\documentclass[11pt, a4paper]{article}

% bfw-style.tex — gemeinsame Style-Definitionen für alle Handouts/Aufgabenhefte
% Voraussetzung: Kompilation mit xelatex (wegen fontspec)

\usepackage[ngerman]{babel}
\usepackage{geometry}
\geometry{a4paper, top=2.2cm, bottom=2.2cm, left=2.3cm, right=2.3cm}

\usepackage{fontspec}
% Fallback-Kette: Source Sans Pro -> Calibri -> Segoe UI -> Latin Modern Sans
% (Latin Modern ist immer mit MiKTeX vorhanden)
\IfFontExistsTF{Source Sans Pro}{%
  \setmainfont{Source Sans Pro}[Ligatures=TeX]%
}{\IfFontExistsTF{Calibri}{%
  \setmainfont{Calibri}[Ligatures=TeX]%
}{\IfFontExistsTF{Segoe UI}{%
  \setmainfont{Segoe UI}[Ligatures=TeX]%
}{%
  \setmainfont{Latin Modern Sans}[Ligatures=TeX]%
}}}
\IfFontExistsTF{Source Code Pro}{%
  \setmonofont{Source Code Pro}[Scale=0.92]%
}{\IfFontExistsTF{Consolas}{%
  \setmonofont{Consolas}[Scale=0.92]%
}{%
  \setmonofont{Latin Modern Mono}[Scale=0.92]%
}}

\usepackage{xcolor}
% Farbpalette: hell, freundlich, modern
\definecolor{bfwPrimary}{HTML}{0EA5A4}    % Petrol/Türkis - Primär
\definecolor{bfwPrimaryDark}{HTML}{0F766E} % dunkler Petrol für Headings
\definecolor{bfwAccent}{HTML}{F59E0B}     % Amber - Akzent
\definecolor{bfwSuccess}{HTML}{10B981}    % Grün - Erfolg / Lösung
\definecolor{bfwWarn}{HTML}{F97316}       % Orange - Achtung
\definecolor{bfwInk}{HTML}{0F172A}        % fast Schwarz
\definecolor{bfwInkSoft}{HTML}{475569}    % gedämpftes Grau
\definecolor{bfwBgSoft}{HTML}{F8FAFC}     % off-white
\definecolor{bfwBgTeal}{HTML}{ECFEFF}     % helles Türkis
\definecolor{bfwBgAmber}{HTML}{FEF3C7}    % helles Gelb
\definecolor{bfwBgRose}{HTML}{FFE4E6}     % helles Rosé für Eselsbrücken

\usepackage[colorlinks=true, linkcolor=bfwPrimaryDark, urlcolor=bfwPrimaryDark]{hyperref}
\usepackage{enumitem}
\setlist[itemize]{leftmargin=*, itemsep=2pt, topsep=2pt}
\setlist[enumerate]{leftmargin=*, itemsep=2pt, topsep=2pt}

\usepackage[most]{tcolorbox}
\usepackage{booktabs}
\usepackage{tikz}
\usetikzlibrary{decorations.pathmorphing, positioning, shapes.geometric, arrows.meta}

% Merkkästchen — wichtige Definition / Kernpunkt
\newtcolorbox{merksatz}[1][]{%
  enhanced, breakable,
  colback=bfwBgTeal,
  colframe=bfwPrimary,
  boxrule=0.6pt,
  arc=4pt,
  left=10pt, right=10pt, top=8pt, bottom=8pt,
  fonttitle=\bfseries\color{bfwPrimaryDark},
  title={\faLightbulb\ Merksatz},
  #1
}

% Eselsbrücke — Mnemoniks
\newtcolorbox{eselsbruecke}[1][]{%
  enhanced, breakable,
  colback=bfwBgRose,
  colframe=bfwAccent,
  boxrule=0.6pt,
  arc=4pt,
  left=10pt, right=10pt, top=8pt, bottom=8pt,
  fonttitle=\bfseries\color{bfwWarn},
  title={Eselsbrücke},
  #1
}

% Beispiel-Box
\newtcolorbox{beispiel}[1][]{%
  enhanced, breakable,
  colback=bfwBgSoft,
  colframe=bfwInkSoft,
  boxrule=0.4pt,
  arc=3pt,
  left=10pt, right=10pt, top=8pt, bottom=8pt,
  fonttitle=\bfseries\color{bfwInk},
  title={Beispiel},
  #1
}

% Lösungs-Box (für Aufgabenhefte)
\newtcolorbox{loesung}[1][]{%
  enhanced, breakable,
  colback=bfwBgAmber!40,
  colframe=bfwSuccess,
  boxrule=0.6pt,
  arc=3pt,
  left=10pt, right=10pt, top=8pt, bottom=8pt,
  fonttitle=\bfseries\color{bfwSuccess!70!black},
  title={Lösung},
  #1
}

% Glossar-Eintrag
\newcommand{\glossareintrag}[2]{%
  \par\noindent\textbf{\color{bfwPrimaryDark}#1} \enspace #2\par\smallskip
}

% Section-Stil — etwas farbiger als Standard
\usepackage{titlesec}
\titleformat{\section}
  {\Large\bfseries\color{bfwPrimaryDark}}
  {\thesection}{0.6em}{}
\titleformat{\subsection}
  {\large\bfseries\color{bfwPrimary}}
  {\thesubsection}{0.5em}{}
\titleformat{\subsubsection}
  {\normalsize\bfseries\color{bfwInk}}
  {\thesubsubsection}{0.4em}{}

% TikZ-Stil "skizzenhaft" — etwas weicher als Rough.js, aber stabil
\tikzset{
  roughbox/.style = {
    draw=bfwPrimaryDark, thick, fill=bfwBgTeal,
    rounded corners=4pt, inner sep=6pt
  },
  roughaccent/.style = {
    draw=bfwAccent, thick, fill=bfwBgAmber,
    rounded corners=4pt, inner sep=6pt
  },
  roughline/.style = {
    draw=bfwInkSoft, thick, -{Stealth[length=2.5mm]}
  }
}

% Bullet-Symbole (bewusst kein FontAwesome — vermeidet Font-Installations-Probleme)
\newcommand{\faLightbulb}{$\bullet$}
\newcommand{\faBookmark}{$\bullet$}

% Header/Footer
\usepackage{fancyhdr}
\pagestyle{fancy}
\fancyhf{}
\renewcommand{\headrulewidth}{0pt}
\fancyfoot[L]{\small\color{bfwInkSoft}\thepage}
\fancyfoot[R]{\small\color{bfwInkSoft} BFW M-064 Beschaffung im Büromanagement}


\title{\vspace*{-1.5cm}\Huge\bfseries\color{bfwPrimaryDark}Aufgabenheft Tag~2\\[0.4em]
  \large\color{bfwInkSoft}\textnormal{Bezugsquellen, Anfragen, Angebotsvergleich --- Übungen im IHK-Klausurformat}}
\author{\textsf{Beschaffung im Büromanagement}\\
\textsf{\small\copyright\ Can Siebert\,\textperiodcentered\,2026}}
\date{Selbstlernmaterial \textperiodcentered{} 12.05.2026}

\begin{document}
\maketitle
\thispagestyle{empty}

\begin{merksatz}[title={\bfseries So nutzen Sie dieses Aufgabenheft}]
Bearbeiten Sie alle Aufgaben zunächst \emph{ohne} in die Lösungen zu schauen.
Beachten Sie die \emph{Operatoren} (nennen, erläutern, beurteilen, begründen,
berechnen, ordnen). Schreiben Sie Antworten in vollständigen Sätzen.
Lösungen ab Seite~\pageref{loesungen}.
\end{merksatz}

\tableofcontents
\vspace{1em}
\hrule

\section{Aufgaben}

\subsection*{Aufgabe~1 --- Bedarfsarten und Brutto-/Netto-Bedarf (10 Punkte)}

\noindent\textbf{\color{bfwAccent}YouTube-Vorbereitung:}\enspace \href{https://www.youtube.com/watch?v=w31S8KZp1Mg}{Bedarfsermittlung einfach erklärt} (\url{https://www.youtube.com/watch?v=w31S8KZp1Mg})

\medskip
Die \emph{Möbelwerk Maier GmbH} fertigt Schreibtische. Für das nächste Quartal sind
3\,000~Schreibtische geplant.

\begin{enumerate}[label=(\alph*)]
  \item \textbf{Ordnen} Sie folgende Materialien der korrekten Bedarfsart zu (Primär-, Sekundär-, Tertiärbedarf): \hfill (4~P.)
    \begin{itemize}[leftmargin=*]
      \item Schreibtische (verkaufsfertig)
      \item Buchenholzplatten
      \item Schmierstoff für die Säge
      \item Tischbeine (zugekaufte Halbfertigprodukte)
    \end{itemize}
  \item \textbf{Erläutern} Sie den Unterschied zwischen \emph{Bruttobedarf} und \emph{Nettobedarf}. \hfill (3~P.)
  \item \textbf{Berechnen} Sie den Nettobedarf für die Buchenholzplatten: Bruttobedarf 6\,500, Lagerbestand 800, bereits bestellt 200, Sicherheitsbestand 300. \hfill (3~P.)
\end{enumerate}

\subsection*{Aufgabe~2 --- ABC-Analyse durchführen (15 Punkte)}

\noindent\textbf{\color{bfwAccent}YouTube-Vorbereitung:}\enspace \href{https://www.youtube.com/watch?v=4u4BPl0aSY8}{ABC-Analyse einfach erklärt mit Berechnung} (\url{https://www.youtube.com/watch?v=4u4BPl0aSY8})

\medskip
Folgende 10 Bürobedarfs-Artikel mit Jahres\-verbrauchswerten liegen vor:

\begin{center}
\begin{tabular}{lr}
\toprule
\textbf{Artikel} & \textbf{Jahresverbrauchswert (€)}\\
\midrule
A1 & 48\,000\\
A2 & 2\,500\\
A3 & 16\,000\\
A4 & 800\\
A5 & 35\,000\\
A6 & 400\\
A7 & 9\,500\\
A8 & 250\\
A9 & 1\,200\\
A10 & 6\,350\\
\bottomrule
\end{tabular}
\end{center}

\begin{enumerate}[label=(\alph*)]
  \item \textbf{Sortieren} Sie die Artikel nach Wert absteigend. \hfill (2~P.)
  \item \textbf{Berechnen} Sie pro Artikel den prozentualen Anteil und den kumulierten Prozentsatz. \hfill (5~P.)
  \item \textbf{Klassifizieren} Sie die Artikel nach ABC (A bis $\approx$\,70\,\%, B bis $\approx$\,90\,\%, C der Rest). \hfill (3~P.)
  \item \textbf{Erläutern} Sie zwei konkrete Konsequenzen der Klassifizierung für die Beschaffungsstrategie. \hfill (5~P.)
\end{enumerate}

\subsection*{Aufgabe~3 --- Anfrage formulieren (10 Punkte)}

\noindent\textbf{\color{bfwAccent}YouTube-Vorbereitung:}\enspace \href{https://www.youtube.com/watch?v=zUf2KHZ_rxk}{Beschaffungsprozess: Anfrage und Bedarfsarten} (\url{https://www.youtube.com/watch?v=zUf2KHZ_rxk})

\medskip
Ihr Betrieb \emph{Schreibwaren-Vertrieb GmbH, Musterstraße 12, 10115 Berlin}, benötigt
50~Aktenordner DIN~A4 (breit, blau, mit Hebelmechanik) und 200~Blatt Druckerpapier
80~g/m².

\begin{enumerate}[label=(\alph*)]
  \item \textbf{Verfassen} Sie eine vollständige schriftliche Anfrage an die \emph{Bürofix AG, Hauptstraße 5, 80331 München}. \hfill (5~P.)
  \item \textbf{Nennen} Sie fünf inhaltliche Bestandteile, die in jeder Anfrage stehen sollten. \hfill (3~P.)
  \item \textbf{Erläutern} Sie: Ist eine Anfrage rechtlich bindend? Begründen Sie Ihre Antwort. \hfill (2~P.)
\end{enumerate}

\subsection*{Aufgabe~4 --- Bezugskalkulation: drei Angebote vergleichen (15 Punkte)}

\noindent\textbf{\color{bfwAccent}YouTube-Vorbereitung:}\enspace \href{https://www.youtube.com/watch?v=OtVKBU2uRDM}{Bezugskalkulation IHK-Prüfung erklärt} (\url{https://www.youtube.com/watch?v=OtVKBU2uRDM})

\medskip
Sie erhalten für 100~Aktenordner drei Angebote:

\begin{center}
\begin{tabular}{lrrrr}
\toprule
\textbf{Angebot} & \textbf{Liste/Stk.} & \textbf{Rabatt} & \textbf{Skonto} & \textbf{Bezugskosten ges.}\\
\midrule
Lieferant A & 4{,}50\,€ & 10\,\% & 2\,\% & 25\,€\\
Lieferant B & 4{,}80\,€ & 12\,\% & 3\,\% & 18\,€\\
Lieferant C & 4{,}20\,€ & 5\,\% & 2\,\% & 30\,€\\
\bottomrule
\end{tabular}
\end{center}

\begin{enumerate}[label=(\alph*)]
  \item \textbf{Berechnen} Sie für jedes Angebot den \textbf{Bezugspreis pro Stück} über das Schema (Liste $-$ Rabatt = Ziel; Ziel $-$ Skonto = Bar; Bar $+$ Bezugskosten/Stk = Bezugspreis). \hfill (9~P.)
  \item \textbf{Beurteilen} Sie: Welcher Lieferant ist quantitativ am günstigsten? \hfill (2~P.)
  \item \textbf{Begründen} Sie, warum der Listenpreis allein kein guter Vergleich ist. \hfill (4~P.)
\end{enumerate}

\subsection*{Aufgabe~5 --- Nutzwertanalyse (Lieferantenbewertung) (12 Punkte)}

\noindent\textbf{\color{bfwAccent}YouTube-Vorbereitung:}\enspace \href{https://www.youtube.com/watch?v=J3l5l--K044}{Nutzwertanalyse einfach erklärt} (\url{https://www.youtube.com/watch?v=J3l5l--K044})

\medskip
Bewerten Sie die drei Lieferanten aus Aufgabe~4 anhand qualitativer Kriterien (Skala
1 = sehr schlecht, 5 = sehr gut):

\begin{center}
\begin{tabular}{lrrrr}
\toprule
\textbf{Kriterium} & \textbf{Gewicht.} & \textbf{A} & \textbf{B} & \textbf{C}\\
\midrule
Lieferzuverlässigkeit & 30\,\% & 4 & 5 & 3\\
Qualität & 25\,\% & 4 & 4 & 5\\
Service / Reklamation & 20\,\% & 3 & 4 & 4\\
Zahlungsziele & 15\,\% & 4 & 3 & 5\\
Nachhaltigkeit & 10\,\% & 3 & 5 & 2\\
\bottomrule
\end{tabular}
\end{center}

\begin{enumerate}[label=(\alph*)]
  \item \textbf{Berechnen} Sie pro Lieferant den gewichteten Nutzwert (Summe aller Punkte $\times$ Gewichtung). \hfill (6~P.)
  \item \textbf{Beurteilen} Sie, welcher Lieferant qualitativ am besten ist. \hfill (2~P.)
  \item \textbf{Vergleichen} Sie das Ergebnis mit der Bezugskalkulation aus Aufgabe~4. Wie würden Sie als Sachbearbeiter/-in endgültig entscheiden? \textbf{Begründen} Sie. \hfill (4~P.)
\end{enumerate}

\subsection*{Aufgabe~6 --- IHK-Wissensfragen (8 Punkte)}

\begin{enumerate}[label=(\alph*)]
  \item \textbf{Erläutern} Sie kurz das Pareto-Prinzip im Kontext der Beschaffung. \hfill (2~P.)
  \item \textbf{Nennen} Sie drei Vorteile eines E-Procurement-Marktplatzes (z.\,B.\ Mercateo) für ein KMU. \hfill (3~P.)
  \item \textbf{Beurteilen} Sie die Aussage: \emph{,,Eine Anfrage und ein Angebot sind dasselbe.''} \hfill (3~P.)
\end{enumerate}

\vspace{1em}
\noindent\textbf{Punkte gesamt: 70}

\newpage

\section{Lösungen}\label{loesungen}

\subsection*{Lösung Aufgabe~1}
\begin{loesung}
\textbf{(a)} Zuordnung:
\begin{itemize}[itemsep=0pt]
  \item Schreibtische (verkaufsfertig): \textbf{Primärbedarf}
  \item Buchenholzplatten: \textbf{Sekundärbedarf} (geht ins Endprodukt ein)
  \item Schmierstoff für die Säge: \textbf{Tertiärbedarf} (Hilfsstoff)
  \item Tischbeine (zugekauft): \textbf{Sekundärbedarf} (geht ins Endprodukt ein)
\end{itemize}

\textbf{(b)} \emph{Bruttobedarf} ist der gesamte Materialbedarf einer Periode ohne Berücksichtigung
von Lagerbestand. \emph{Nettobedarf} ist die Menge, die tatsächlich neu beschafft werden muss
($=$ Brutto $-$ Lager $-$ offene Bestellungen $+$ Sicherheitsbestand).

\textbf{(c)} Nettobedarf Buchenholzplatten:
$6\,500 - 800 - 200 + 300 = \textbf{5\,800 Stück}$.
\end{loesung}

\subsection*{Lösung Aufgabe~2}
\begin{loesung}
\textbf{(a)+(b)+(c)} Sortiert nach Wert absteigend mit Anteilen:
Gesamtwert = 125\,200~€.

\begin{center}
\begin{tabular}{lrrrl}
\toprule
\textbf{Artikel} & \textbf{Wert (€)} & \textbf{\%} & \textbf{Kumul.} & \textbf{Klasse}\\
\midrule
A1 & 48\,000 & 38{,}3\,\% & 38{,}3\,\% & A\\
A5 & 35\,000 & 27{,}9\,\% & 66{,}2\,\% & A\\
A3 & 16\,000 & 12{,}8\,\% & 79{,}0\,\% & B\\
A7 & 9\,500 & 7{,}6\,\% & 86{,}6\,\% & B\\
A10 & 6\,350 & 5{,}1\,\% & 91{,}7\,\% & B\\
A2 & 2\,500 & 2{,}0\,\% & 93{,}7\,\% & C\\
A9 & 1\,200 & 1{,}0\,\% & 94{,}6\,\% & C\\
A4 & 800 & 0{,}6\,\% & 95{,}3\,\% & C\\
A6 & 400 & 0{,}3\,\% & 95{,}6\,\% & C\\
A8 & 250 & 0{,}2\,\% & 95{,}8\,\% & C\\
\bottomrule
\end{tabular}
\end{center}

(Hinweis: kleine Rundungs­ungenauigkeiten ergeben Schluss bei 95{,}8\,\% statt 100\,\% --- das ist ok.)

\textbf{(d)} Konsequenzen für die Beschaffungs­strategie:
\begin{itemize}[itemsep=0pt]
  \item \emph{A-Teile (A1, A5)}: Rahmenverträge, häufige Bestellzyklen, kleine Sicherheits­bestände, intensive Lieferantenpflege.
  \item \emph{B-Teile (A3, A7, A10)}: Periodische Bestellung (z.\,B.\ monatlich), mittlere Sicherheits­bestände, regelmäßige Audits.
  \item \emph{C-Teile (Rest)}: Vorratsbeschaffung mit großzügigen Sicherheits­beständen, Standardprozess, kein Verhandlungsaufwand.
\end{itemize}
\end{loesung}

\subsection*{Lösung Aufgabe~3}
\begin{loesung}
\textbf{(a)} Musterbeispiel einer Anfrage (Auszug, vereinfacht):

\begin{quote}\small
Schreibwaren-Vertrieb GmbH \\
Musterstraße 12 \\
10115 Berlin \\

\medskip
Bürofix AG \\
Hauptstraße 5 \\
80331 München \\

\medskip
Berlin, 12.05.2026 \\

\medskip
\textbf{Anfrage Aktenordner und Druckerpapier}

\medskip
Sehr geehrte Damen und Herren, \\

\medskip
für unseren Bürobedarf benötigen wir folgende Artikel und bitten Sie um ein Angebot:

\begin{itemize}[itemsep=0pt]
  \item 50~Stück Aktenordner DIN~A4, breit, blau, mit Hebelmechanik
  \item 200~Blatt Druckerpapier DIN~A4, 80~g/m², weiß
\end{itemize}

\smallskip
Bitte teilen Sie uns Ihre Liefer- und Zahlungs­bedingungen, Liefer­zeit, Verpackungs- sowie
Bezugskosten mit. Wir bitten um Ihr Angebot bis zum \textbf{19.05.2026}.

\medskip
Mit freundlichen Grüßen \\
Ihr Name --- Sachbearbeitung Einkauf
\end{quote}

\textbf{(b)} Fünf Bestandteile: Artikel­bezeichnung, Menge, gewünschter Liefertermin,
Lieferort, Frage nach Liefer- und Zahlungsbedingungen (alternativ: Verpackung, Garantie, Angebotsfrist).

\textbf{(c)} Eine Anfrage ist \textbf{rechtlich nicht bindend}. Sie ist eine
\emph{invitatio ad offerendum} --- eine unverbindliche Aufforderung an den Lieferanten,
ein Angebot abzugeben. Erst dessen Angebot ist nach \S~145 BGB ein bindender Vertragsantrag.
\end{loesung}

\subsection*{Lösung Aufgabe~4}
\begin{loesung}
\textbf{(a)} Bezugskalkulation für jedes Angebot (alle Werte pro Stück):

\begin{center}
\begin{tabular}{lrrr}
\toprule
\textbf{Position} & \textbf{Lief.\,A} & \textbf{Lief.\,B} & \textbf{Lief.\,C}\\
\midrule
Listenpreis & 4{,}50\,€ & 4{,}80\,€ & 4{,}20\,€\\
$-$ Liefererrabatt & 0{,}45 (10\,\%) & 0{,}58 (12\,\%) & 0{,}21 (5\,\%)\\
$=$ Zieleinkaufspreis & 4{,}05\,€ & 4{,}22\,€ & 3{,}99\,€\\
$-$ Liefererskonto & 0{,}08 (2\,\%) & 0{,}13 (3\,\%) & 0{,}08 (2\,\%)\\
$=$ Bareinkaufspreis & 3{,}97\,€ & 4{,}09\,€ & 3{,}91\,€\\
$+$ Bezugskosten / Stück & 0{,}25 & 0{,}18 & 0{,}30\\
\midrule
\textbf{Bezugspreis / Stück} & \textbf{4{,}22\,€} & \textbf{4{,}27\,€} & \textbf{4{,}21\,€}\\
\bottomrule
\end{tabular}
\end{center}

(Bezugskosten/Stück = Bezugskosten gesamt / 100~Stück.)

\textbf{(b)} Quantitativ ist Lieferant \textbf{C} mit 4{,}21\,€ am günstigsten,
knapp vor A (4{,}22\,€). Lieferant B ist mit 4{,}27\,€ am teuersten --- trotz höchstem
Rabatt, weil sein Listenpreis am höchsten war.

\textbf{(c)} Der Listenpreis allein ist kein guter Vergleich, weil er weder Rabatte noch
Skonto noch Bezugskosten enthält. Lieferant B mit dem höchsten Rabatt (12\,\%) wirkt
zunächst attraktiv, ist aber wegen seines hohen Listenpreises tatsächlich am teuersten.
Nur der Bezugspreis (Einstandspreis) erlaubt einen fairen Vergleich.
\end{loesung}

\subsection*{Lösung Aufgabe~5}
\begin{loesung}
\textbf{(a)} Gewichteter Nutzwert pro Lieferant:

\begin{center}
\begin{tabular}{lrrrr}
\toprule
\textbf{Kriterium} & \textbf{Gew.} & \textbf{A} & \textbf{B} & \textbf{C}\\
\midrule
Lieferzuverlässigkeit & 0{,}30 & 4 $\to$ 1{,}20 & 5 $\to$ 1{,}50 & 3 $\to$ 0{,}90\\
Qualität & 0{,}25 & 4 $\to$ 1{,}00 & 4 $\to$ 1{,}00 & 5 $\to$ 1{,}25\\
Service & 0{,}20 & 3 $\to$ 0{,}60 & 4 $\to$ 0{,}80 & 4 $\to$ 0{,}80\\
Zahlungsziele & 0{,}15 & 4 $\to$ 0{,}60 & 3 $\to$ 0{,}45 & 5 $\to$ 0{,}75\\
Nachhaltigkeit & 0{,}10 & 3 $\to$ 0{,}30 & 5 $\to$ 0{,}50 & 2 $\to$ 0{,}20\\
\midrule
\textbf{Nutzwert gesamt} & 1{,}00 & \textbf{3{,}70} & \textbf{4{,}25} & \textbf{3{,}90}\\
\bottomrule
\end{tabular}
\end{center}

\textbf{(b)} Qualitativ gewinnt Lieferant \textbf{B} klar mit Nutzwert 4{,}25, vor C (3{,}90)
und A (3{,}70).

\textbf{(c)} \emph{Synthese}: Quantitativ vorne C (4{,}21\,€/Stück), qualitativ vorne B
(Nutzwert 4{,}25). Bei nur 0{,}06\,€ Aufpreis/Stück gegenüber C bietet B deutlich höhere
Lieferzuverlässigkeit (5 vs.\ 3) und bessere Nachhaltigkeit (5 vs.\ 2). Bei 100~Stück
sind das 6,00\,€ Mehrkosten gegen ein deutlich geringeres Reklamations- und Lieferausfall­
risiko --- ein guter Tausch. \textbf{Empfehlung: Lieferant B} beauftragen, mit der Begründung
„Mehrkosten von 6\,€ rechtfertigen sich durch deutlich höhere Lieferzuverlässigkeit''.
\end{loesung}

\subsection*{Lösung Aufgabe~6}
\begin{loesung}
\textbf{(a)} Pareto-Prinzip: 20\,\% der Ursachen verursachen 80\,\% der Wirkung. In der
Beschaffung: Etwa 20\,\% der Artikel machen 80\,\% des wirtschaftlichen Werts aus
(A-Teile). Konsequenz: Ressourcen auf die wenigen wichtigen Artikel konzentrieren,
nicht gleichmäßig verteilen.

\textbf{(b)} Drei Vorteile eines Marktplatzes (z.\,B.\ Mercateo) für ein KMU:
\begin{itemize}[itemsep=0pt]
  \item Keine eigene IT-Investition (Plattform wird vom Marktplatz-Betreiber bereitgestellt).
  \item Breite Lieferantenauswahl an einem Ort --- vergleichen ohne mehrere Webshops.
  \item Einheitlicher Bestellprozess und konsolidierte Rechnung über den Marktplatz.
\end{itemize}

\textbf{(c)} Die Aussage ist falsch. Eine \emph{Anfrage} ist eine rechtlich \textbf{unverbindliche}
Aufforderung an den Lieferanten, ein Angebot abzugeben (\emph{invitatio ad offerendum}).
Ein \emph{Angebot} ist hingegen ein \textbf{bindender Vertragsantrag} (§~145 BGB), der bei
Annahme des Käufers zum Vertragsschluss führt. Anfrage und Angebot sind also rechtlich
deutlich unterschiedliche Willenserklärungen.
\end{loesung}

\end{document}
